\chapter{Introduction}

Replication is necessary for fault tolerance, but in a high-throughput
transaction system it can easily become the limiting stage. A database can use
many CPU cores to execute transactions, yet still lose scalability if every
completed transaction must wait behind one slow replication or replay step.
This is the setting of \Mako{}, a geo-replicated transactional key-value store
whose original replication backend was Multi-Paxos \cite{mako-osdi25}.

\Mako{} was designed around a pipeline. A leader-side process accepts
transaction work and uses worker threads to execute many transactions in
parallel. The effects of those transactions are recorded as transaction log
records. A replication backend orders the records, and follower replicas replay
the committed records into their local database state. Chapter 2 defines these
terms in detail; the important point for the introduction is that \Mako{}'s
performance depends on all three stages staying healthy: leader execution,
replication commit, and follower replay.

This thesis studies a practical systems question:

\begin{center}
\emph{Can Raft replace Paxos in Mako while preserving high throughput?}
\end{center}

Raft is attractive because it packages leader election, log replication, and
leader changes into a protocol that is easier to understand and operate than
Paxos \cite{raft-atc14}. Paxos and production Multi-Paxos systems provide the
baseline consensus background for this comparison
\cite{paxos-made-simple,paxos-made-live}. Raft's understandability advantage
is useful only if Raft can fit \Mako{}'s performance shape. \Mako{} is not a
small replicated service with one command stream. Its replication path
interacts with worker threads, partitioned transaction lanes, follower replay,
leader placement, and benchmark measurements that report leader-side
throughput.

The hard part is that consensus correctness and end-to-end system performance
are different properties. A Raft backend can order commands correctly and
still be a poor replacement if it funnels too much work through one narrow
apply path, elects a leader that does not match \Mako{}'s transaction
submission path, or lets followers replay committed work too slowly. The
measured throughput must represent real replicated progress, not only records
accepted at the leader.

The central argument of this thesis is that Raft can replace \Mako{}'s Paxos
backend without leaving Paxos's performance class, but only when the
integration preserves the pipeline around the consensus protocol. Multi-Raft
keeps the original Paxos-like parallel structure by using separate Raft groups
for separate transaction lanes. Single-Raft reduces the replication topology by
routing all partition commands through one consolidated Raft log, with one
replica of that group running in each process, but then follower replay becomes
the bottleneck that must be engineered explicitly. ReplayPool repairs that
bottleneck by moving heavy follower replay work out of the narrow Raft apply
path. Preferred leader election is used by the Raft integration more
generally: both Multi-Raft and Single-Raft must align Raft leadership with the
process where this \Mako{} benchmark submits and first executes transactions.

\section{Why a Straight Protocol Swap Is Not Enough}

It is tempting to describe this work as replacing one consensus protocol with
another. That description is too small. \Paxos{} and \Raft{} both provide a
way for replicas to agree on an ordered log, but the protocol is only one part
of \Mako{}'s execution path. The surrounding system decides where transaction
work is submitted, how many worker lanes produce log records, how committed
records are delivered back into \Mako{}, and whether followers can apply those
records fast enough.

This distinction matters because a benchmark can otherwise be misleading. If
the leader process accepts and commits records quickly, the reported
leader-side throughput can look strong. But if follower replicas replay those
records much more slowly, the system is accumulating replicated work faster
than followers can absorb it. A thesis about replacing \Paxos{} in \Mako{}
therefore needs more than a throughput curve. It needs evidence about topology,
leader placement, replay progress, CPU use, and persistence work.

\section{Design Overview}

The first design point is Multi-Raft. Multi-Raft is the conservative
replacement because it preserves the original shape of the Paxos path. Each
worker/partition lane has its own Raft group, so independent streams of
transaction records can still be ordered independently. Multi-Raft answers the
baseline question: if Raft is given a Paxos-like topology, can it stay in the
same performance class?

The second design point is Single-Raft. Single-Raft asks whether the
replication topology can be simplified to one consolidated Raft log for all
partition commands, with one replica of that log running in each process. This
reduces replication structure, but it also concentrates committed work into
one ordered Raft path. The apply path is the code path that handles a committed
Raft entry after consensus has ordered it. Replay work is the heavier
follower-side step of applying that committed transaction record to local
database state. Single-Raft is useful only if consolidation does not turn
apply or follower replay into a hidden serial bottleneck.

The third design point is preferred leader election. \Mako{}'s transaction
submission path expects a particular process to be the leader-side producer of
transaction records. Standard Raft can elect any eligible replica. Preferred
leader election biases leadership toward the process that owns the Mako
submission path while preserving normal Raft safety and failover rules. It is a
full integration mechanism rather than a cosmetic preference: startup
placement makes the normal Raft leader match Mako's submission process, backup
replicas can still lead if that process is unavailable, and leadership can move
back only after the preferred replica has caught up. This placement mechanism
is needed to align this leader-oriented Mako integration and evaluation with
the expected submission process; it is not a new safety condition for Raft. It
applies to the Raft integration as a whole, not only to Single-Raft. Chapter 4
explains this mechanism in depth.

The fourth design point is ReplayPool. ReplayPool makes Single-Raft viable by
keeping Raft apply lightweight. The apply path stages committed payloads and
dispatches replay work; parallel replay workers perform the heavier
follower-side database application. In this thesis, ``1:1'' replay
configuration means the number of replay workers matches the number of Mako
worker threads. It does not mean one Raft group per worker.

\section{Contributions}

This work makes five contributions.

First, it adds a runtime-selectable Raft backend for Mako. The transaction
layer submits transaction logs through a shared replication interface, allowing
Paxos, Multi-Raft, and Single-Raft to run under the same workload, benchmark
harness, and metric extraction path. This makes the comparison a systems
comparison rather than a collection of unrelated experiments.

Second, it implements preferred leader election for Mako. The mechanism biases
Raft leadership toward the process that owns leader-side transaction
submission for both Multi-Raft and Single-Raft, while preserving normal Raft
voting, log matching, and failover. It also supports the failure case that
matters for a leader-placement policy: if the preferred replica is unavailable,
a backup replica can lead, and leadership is moved back only after the preferred
replica catches up. This is a placement mechanism around Raft, not a change to
Raft's safety rules.

Third, it evaluates Multi-Raft and Single-Raft as two different Raft
integration strategies. Multi-Raft preserves Paxos-like partition-level
parallelism. Single-Raft consolidates partition commands into one process-wide
Raft log and tests whether a simpler topology can still support Mako's
execute-replicate-replay pipeline.

Fourth, it adds ReplayPool offload for Single-Raft. ReplayPool keeps Raft
apply lightweight and moves heavy follower replay work into parallel Mako
replay workers. This makes follower replay throughput part of the design
evidence, not an afterthought hidden behind leader-side throughput.

Fifth, it adds source-tagged simulated-disk instrumentation. The simulated-disk
counters report bytes, writes, and normalized storage work so persistence
results are tied to measured work rather than only to throughput curves. The
simulated storage model is intentional: real SSD and cloud-storage behavior
varies by device, service, and configuration. This thesis therefore evaluates
representative NVMe and Cloud-SSD conditions with controlled bandwidth,
latency, and queueing assumptions, then checks whether the backends performed
comparable measured storage work.

\section{Evidence Preview}

The evaluation follows the same story. In the in-memory experiments, Multi-Raft
tracks the original Paxos backend closely, and Single-Raft with ReplayPool
stays in the same throughput class while keeping follower replay much closer
to leader-side committed throughput. Single-Raft without ReplayPool is an
important negative result: it can report high leader-side throughput while
followers replay much more slowly. Under simulated Cloud-SSD persistence, the
backends flatten near the same throughput, and the byte/write counters show
comparable modeled storage work. These results support the thesis claim that
Raft is viable in \Mako{} when topology, leader placement, replay parallelism,
and persistence accounting are treated as part of the replication design.

The persistence experiments include both simulated NVMe and simulated
Cloud-SSD settings. The NVMe setting checks whether a faster storage model
leaves the replication backends close to their in-memory behavior. The
Cloud-SSD setting applies a slower, more restrictive storage model and is the
case where the throughput curves flatten most clearly. In both cases, the
interpretation depends on the measured bytes and writes, not on the shape of a
throughput curve alone.

The evaluation is intentionally organized around claims rather than around
implementation files. The no-disk scalability experiments test whether Raft
can stay in \Paxos{}'s performance class when storage is removed from the
critical path. The worker CPU measurements check whether \Mako{} workers are
still doing useful transaction work. The replay measurements check whether
followers keep up with committed records. The simulated-disk measurements
attach byte and write-count evidence to the persistence results so storage
claims are not based on throughput shape alone.

The preferred-leader mechanism is validated separately from the headline
throughput graphs. Dedicated tests check startup placement, log replication
through the preferred leader, ordinary backup leadership while the preferred
replica is absent, and safe return to the preferred replica after it catches
up. A fresh Single-Raft and Multi-Raft sweep over 1 through 11 worker threads
was also run as a no-regression check. Those checks support the leader-placement
claim without replacing the frozen evaluation results used in Chapter 8.

\section{Thesis Roadmap}

The rest of the thesis follows the structure of that argument. Chapter 2
describes the Mako pipeline and original Paxos path. Chapter 3 explains the
Raft backend and the topology choices. Chapter 4 presents preferred leader
election. Chapter 5 explains the Single-Raft replay bottleneck and ReplayPool.
Chapter 6 describes persistence and simulated disk. Chapter 7 summarizes
validation. Chapter 8 evaluates the implementation. Chapter 9 discusses
limitations and tradeoffs. Chapter 10 relates the work to prior systems, and
Chapter 11 concludes.
